\documentclass{../../slides-style}

\slidetitle[Введение]{Операционные системы, практика}{05.09.2026}

\author[]{}

\begin{document}

    \begin{frame}[plain]
        \titlepage
    \end{frame}

    \section{Формальные вопросы}

    \begin{frame}{Формальные вопросы}
        \begin{outline}
            \1 Занятия по субботам раз в две недели в 3389
                \2 Две пары практики сразу, чтобы успеть сделать что-то содержательное
                \2 Будем работать прежде всего на парах, домашек не будет (если только доделать что-то)
                \2 Работа в основном командная
            \1 Берите с собой ноуты с qemu
            \1 Курс на HwProj: \url{https://hwproj.ru/courses/50095}
            \1 Контакты преподавателей:
                \2 Никита Юрьевич Ловягин: n.lovyagin@spbu.ru
                \2 Михаил Александрович Серов: m.serov@spbu.ru
                \2 Георгий Константинович Сичкар: georgy.sichkar@gmail.com
                \2 Антон Владимирович Черноусов: anton.chernousov@spbu.ru
                \2 Юрий Викторович Литвинов: y.litvinov@spbu.ru
                \2 Яков Александрович Кириленко: y.kirilenko@spbu.ru
                \2 Сергей Владимирович Погожев: s.pogozhev@spbu.ru
        \end{outline}
    \end{frame}

    \begin{frame}{Критерии оценивания}
        \begin{outline}
            \1 Всего примерно 7 заданий, каждое оценивается из 10 баллов
                \2 Можно один раз пересдать, чтобы улучшить балл
                \2 -3 невосстановимых за пропуск дедлайна
                \2 Делаем в основном на паре, но если пропустили, можно досдать
            \1 Баллы за домашние работы переводятся в шкалу 0-40
            \1 Будут две \enquote{контрольные} --- индивидуальные опросы и простые задачи на листочке
                \2 так что даже в командах надо понимать \textit{всё}, что делается
            \1 Могут быть летучки с теоретическими вопросами на 5 минут, на любой паре
            \1 В конце --- теоретический зачёт по шкале 0-60
        \end{outline}
    \end{frame}

    \begin{frame}{Что будет в курсе}
        \begin{outline}
            \1 Базовая работа в xv6 --- установка/запуск, написание простых приложений с системными вызовами, добавим свой системный вызов
            \1 Работа с таблицей процессов
            \1 Многопоточность и синхронизация, мьютексы
            \1 Трассировка событий ядра, dmesg
            \1 Межпроцессное взаимодействие
            \1 Написание драйверов
            \1 Файловые системы --- чтение образа ext2
            \1 Вполне возможно, что-то ещё
        \end{outline}
        Практики могут немного обгонять лекции!
    \end{frame}

    \section{xv6}

    \begin{frame}{Краткий обзор xv6}
        \begin{outline}
            \1 Учебная операционная система для курса MIT по ОС
                \2 \url{https://pdos.csail.mit.edu/6.1810/2026/index.html}
            \1 POSIX-like
                \2 Не поддерживает много чего из POSIX, зато простая
                \2 \enquote{Почти Linux}
            \1 Важные отличия от обычных ОС:
                \2 Один поток на процесс
                \2 Не линкуется со стандартной библиотекой C (glibc), есть своя сильно урезанная версия (ulib)
                \2 Собираем образ прямо с прикладной программой (как во встроенных системах)
                \2 Заточена под 64-битный RISC-V
            \1 Репо: \url{https://github.com/mit-pdos/xv6-riscv}
            \1 Инструкция по сборке и установке: \url{https://pdos.csail.mit.edu/6.1810/2026/tools.html}
        \end{outline}
    \end{frame}

    \begin{frame}{Структура исходников}
        \begin{outline}
            \1 kernel --- монолитное ядро системы
                \2 entry.S --- начальный загрузчик (на ассемблере RISC-V)
                \2 main.c --- собственно инициализация системы
                \2 proc.c --- процессы и планировщик
                \2 trap.c --- обработка прерываний
                \2 syscall.c --- диспетчер системных вызовов
                \2 console.c --- код для работы с консолью
                \2 pipe.c --- код работы с \textit{пайпами}
                \2 ...
            \1 user --- программы пространства пользователя
                \2 sh.c --- командная оболочка
                \2 cat.c, echo.c, grep.c, kill.c, wc.c --- базовые утилиты
                    \3 не стесняйтесь почитать их код, они очень простые!
                \2 ls.c, mkdir.c, ln.c, rm.c --- утилиты для работы с файлами
                \2 user.h --- заголовочный файл для системных вызовов, ulib и printf 
                    \3 его подключать в свои программы
            \1 mkfs --- утилита, собирающая загрузочный образ
        \end{outline}
    \end{frame}

    \begin{frame}{Системные вызовы}
        \begin{outline}
            \1 int fork() --- создаёт новый процесс, копию текущего
            \1 int exit(int status) --- заканчивает текущий процесс
            \1 int wait(int *status) --- блокирует текущий процесс, пока не закончится дочерний
            \1 int pause(int n) --- блокирует текущий процесс на n тиков
            \1 int getpid() --- возвращает ID текущего процесса
            \1 int kill(int pid) --- завершить указанный процесс
            \1 int dup(int fd) --- создать новый файловый дескриптор для fd
            \1 int pipe(int p[]) --- создать пайп, возвращает в p два дескриптора
            \1 \dots
        \end{outline}
    \end{frame}

    \begin{frame}[fragile]{Процессы}
        \begin{outline}
            \1 У каждого процесса свой pid, можно узнать с помощью \mintinline{c}|pid_t getpid(void);|
            \1 Программно создать процесс --- \mintinline{c}|pid_t fork(void);|
                \2 Возвращает 0 для дочернего процесса и pid дочернего процесса для родительского, -1, если ошибка
                \2 Дочерний процесс является копией родительского, данные меняются \emph{независимо}
            \1 Завершение процесса --- \mintinline{c}|void exit(int status);|
                \2 return из main вызывает exit
            \1 Дождаться завершения дочернего процесса --- \mintinline{c}|pid_t wait(int *wstatus);|
                \2 Если дочерних процессов несколько, ждём первого завершившегося
            \1 Заместить текущий процесс бинарником с диска --- \mintinline{c}|int exec(char *file, char *argv[])|
                \2 Обычно используется по схеме fork --- настройка потоков ввода-вывода --- exec
        \end{outline}
    \end{frame}

    \begin{frame}[fragile]{Ввод-вывод}
        \begin{outline}
            \1 С каждым процессом ассоциирована \emph{таблица дескрипторов} открытых файлов
                \2 Аналог FILE* из glibc, но в xv6 в пользовательских программах это просто число
            \1 Дескриптор с индексом 0 --- stdin, 1 --- stdout, 2 --- stderr
            \1 Их можно подменять: \mintinline{c}|close(0); open("input.txt", O_RDONLY);|
                \2 open гарантированно выберет минимальный свободный индекс
                \2 в процессе всегда один поток, так что никто не \enquote{украдёт} ячейку под дескриптор
        \end{outline}
    \end{frame}

    \begin{frame}[fragile]{Пример: пайп}
        \begin{footnotesize}
            \begin{minted}{c}
int p[2];
char *argv[2];
argv[0] = "wc";
argv[1] = 0;
pipe(p);
if (fork() == 0) {
    close(0); // Закрываем stdin
    dup(p[0]); // Ячейка 0 свободна, копия выходного конца пайпа "упадёт" туда
    close(p[0]); // Закрываем исходный дескриптор
    close(p[1]); // Закрываем входной конец пайпа
    exec("/bin/wc", argv); // Стартуем wc с подменённым stdin
} else {
    close(p[0]); // Закрываем выходной конец пайпа
    write(p[1], "hello world\n", 12); // Пишем в пайп
    close(p[1]); // Закрываем входной конец (иначе wc не закончит работу)
}
            \end{minted}
        \end{footnotesize}
    \end{frame}

    \begin{frame}[fragile]{Как добавить свою программу}
        \begin{outline}
            \1 Добавить файл .c в папку user (проще всего скопировать что-то типа cat)
            \1 Не забыть подключить user.h, ещё могут быть полезны kernel/types.h и kernel/stat.h
            \1 Реализовать там содержательную логику
                \2 система сборки не любит многофайловые утилиты, так что лучше всё в одном файле
            \1 Добавить в Makefile:
            \begin{minted}{make}
UPROGS=\ ...
    $U/_filename
            \end{minted}
            \1 Собрать и запустить систему (make qemu)
            \1 Вызвать вашу программу из шелла как обычно
        \end{outline}
    \end{frame}

    \begin{frame}[fragile]{Как добавить системный вызов}
        \begin{footnotesize}
            \begin{outline}
                \1 Добавить в файл kernel/syscall.h константу с номером системного вызова, начинающуюся с префикса \mintinline{text}|SYS_|
                \1 Добавить прототип функции ядра обработки системного вызова в файл kernel/syscall.c
                    \2 \footnotesize Функция не должна иметь аргументов и должна иметь возвращаемое значение типа uint64
                \1 В том же файле добавить указатель на эту функцию в таблицу системных вызовов (syscalls)
                \1 Создать в kernel отдельный .c с префиксом \mintinline{text}|sys_|, где и реализовать содержательную логику
                \1 Добавить прототип системного вызова в user/user.h
                \1 Добавить функцию в user/usys.pl, чтобы сгенерировалась машинерия вызова из пространства пользователя
                \1 Добавить в Makefile в OBJS: \mintinline{text}|$K/sys<ваш системный вызов>.o|
                \1 Принимать параметры внутри системного вызова надо через функцию argint
                    \2 \footnotesize У ядра свой стек, обычная передача параметров не работает
                    \2 \footnotesize Например, заголовку \mintinline{text}|int add (int, int);| в user.h соответствует \mintinline{text}|uint64 sys_add(void)| в ядре
            \end{outline}
        \end{footnotesize}
    \end{frame}

    \section{Задачи}

    \begin{frame}{Задача 1}
        \framesubtitle{На 3 балла}
        Написать программу, выполняющую следующие действия в xv6.
        \begin{outline}
            \1 Создается дочерний процесс, который приостанавливается (pause) 5-15 секунд и завершается с кодом возврата 1 (exit).
            \1 Родительский процесс выводит свой идентификатор (getpid) и идентификатор дочернего процесса, затем
                \2 (a) просто ожидает завершения дочернего процесса (wait), после чего выводит идентификатор завершенного процесса и его код возврата, затем сам завершается (exit)
                \2 (b) сначала посылает дочернему процессу сигнал kill, после чего ожидает его завершения, выводит идентификатор завершенного процесса и код возврата, затем сам завершается
            \1 В варианте (a) протестировать запуск приложения в виде фонового процесса оболочки и отправку сигнала дочернему процессу с помощью команды kill, не дожидаясь его завершения по тайм-ауту.
        \end{outline}
    \end{frame}

    \begin{frame}[fragile]{Задача 2}
        \framesubtitle{На 3 балла}
        Добавить в xv6 системный вызов, вычисляющий сумму двух целых чисел (своих аргументов). 
        \begin{outline}
            \1 Для доступа к аргументам системного вызова использовать соответствующие функции (argint)
            \1 Вернуть сумму чисел как результат системного вызова
            \1 Добавить отладочный вывод
            \1 Протестировать работу с помощью утилиты пользовательского пространства
        \end{outline}
    \end{frame}

    \begin{frame}[fragile]{Задача 3* (1)}
        \framesubtitle{На 4 балла}
        Написать программу, которая создает дочерний процесс и передает ему через заранее созданный анонимный канал (pipe) все аргументы командной строки с помощью write, завершая их символом перевода строки \mintinline{c}|'\n'|, закрывает канал по завершении записи. 
        
        Дочерний процесс замещает стандартный ввод (0) выходным концом канала (\mintinline{c}|pipefd[0]|) с помощью dup:

        \begin{minted}{c}
close(0);
dup(pipefd[0]);
        \end{minted}

        после чего канал может быть закрыт
        
        \begin{minted}{c}
close(pipefd[0]);
        \end{minted}
    \end{frame}

        \begin{frame}[fragile]{Задача 3* (2)}
        \framesubtitle{На 4 балла}
        Далее дочерний процесс замещается приложением wc

        \begin{minted}{c}
char *argv[] = {"/wc", 0};
exec("/wc", argv);
        \end{minted}

        Программа wc должна вывести подсчет строк (количества аргументов командной строки), слов и байтов в них. Родительский процесс ожидает завершение дочернего процесса (wait) и завершается сам (exit).
    \end{frame}

    \begin{frame}{Задача 4}
        \framesubtitle{На 3 балла, бонусная}
        Реализовать в \enquote{настоящей} ОС (POSIX) программу, которая создает дочерний процесс и передает ему через заранее созданный анонимный канал (pipe) все аргументы командной строки с помощью write, завершая их символом перевода строки \mintinline{c}|'\n'|, закрывает канал по завершении записи. Дочерний процесс читает данные из канала с помощью read и выводит их на стандартный вывод.
    \end{frame}

    \begin{frame}{Обратите внимание! (1)}
        \begin{outline}
            \1 Вызов pause в xv6 считает \enquote{ticks} (примерно 1/10 секунды, вычисляется по инструкциям процессора)
                \2 подсчета \enquote{настоящего} времени нет
            \1 Закрывайте используемые концы канала и иные файлы после работы
                \2 ОС их закроет при завершении процесса, но лучше самим
            \1 Закрытие неиспользуемых концов канала обязательно, т.к. иначе может возникнуть взаимоблокировка
            \1 Проверяйте на ошибки все используемые функции и системные вызовы!
        \end{outline}
    \end{frame}

    \begin{frame}{Обратите внимание! (2)}
        \begin{outline}
            \1 Вызов exec может завершиться с ошибкой, тогда исполняется следующий за ним код процесса
                \2 иначе исполняется запущенная программа, про исходный процесс \enquote{забывают}
            \1 Оболочка xv6 не поддерживает экранирование пробелов и других символов (кавычки и т.п.), поэтому wc будет возвращать равное количество слов и строк
                \2 можно вызывать wc из другой с помощью exec, и передать аргументы с пробелами или переводом строки
                \2 искусственное добавление \mintinline{c}|'\n'| к каждому аргументу в родительском процессе не гарантирует, что wc посчитает число аргументов исходного процесса, тем более в реальной системе
        \end{outline}
    \end{frame}

    \begin{frame}{Обратите внимание! (3)}
        \begin{outline}
            \1 Оболочка xv6 по каким-то причинам ограничивает количество аргументов приложения сильнее, чем ядро (10 против 32). Это можно пропатчить.
            \1 Формально, успехом write считается любое положительное число выведенных байтов
                \2 write может вывести меньше данных, чем указано, и это не будет ошибкой --- надо следить за реально выведенными байтами
            \1 аналогичная ситуация с read: реальное количество прочитанных байтов может быть меньше размера буфера 
                \2 чтение нужно проводить до достижения конца файла (0) или до ошибки (отрицательное значение)
        \end{outline}
    \end{frame}

\end{document}

